\documentclass[11pt,a4paper]{article}

%% ===== Core Packages =====
\usepackage{fontspec}
\usepackage{polyglossia}
\setmainlanguage{english}
\usepackage{xeCJK}
\usepackage{amsmath,amssymb,amsthm}
\usepackage{geometry}
\usepackage{graphicx}
\usepackage{xcolor}
\usepackage{titlesec}
\usepackage{fancyhdr}
\usepackage{enumitem}
\usepackage{array,booktabs,longtable,multirow}
\usepackage{hyperref}
\usepackage{tocloft}
\usepackage{indentfirst}
\usepackage{tikz}

%% ===== Page Geometry =====
\geometry{
  paperwidth=210mm,paperheight=297mm,
  left=22mm,right=22mm,top=28mm,bottom=28mm,
  headheight=14pt,footskip=30pt
}

%% ===== Font Configuration =====
\setmainfont{Linux Libertine O}
\setsansfont{Linux Biolinum O}
\setmonofont{DejaVu Sans Mono}[Scale=MatchLowercase]

%% ===== CJK Font Setup =====
\setCJKmainfont[Path=/mnt/agents/fonts/noto-sc-extracted/]{NotoSansCJKsc-Regular.otf}
\setCJKsansfont[Path=/mnt/agents/fonts/noto-sc-extracted/]{NotoSansCJKsc-Regular.otf}
\setCJKmonofont[Path=/mnt/agents/fonts/noto-sc-extracted/]{NotoSansCJKsc-Regular.otf}

%% ===== Colors =====
\definecolor{titlecolor}{RGB}{45,55,72}
\definecolor{sectioncolor}{RGB}{55,65,81}
\definecolor{notecolor}{RGB}{120,53,15}
\definecolor{rulecolor}{RGB}{180,160,140}
\definecolor{coverblue}{RGB}{35,55,85}

%% ===== Section Formatting =====
\titleformat{\section}
  {\normalfont\Large\bfseries\color{sectioncolor}}%
  {\thesection}{1em}{}
\titleformat{\subsection}
  {\normalfont\large\bfseries\color{sectioncolor}}
  {\thesubsection}{1em}{}
\titleformat{\subsubsection}
  {\normalfont\normalsize\bfseries\color{sectioncolor}}
  {\thesubsubsection}{1em}{}

%% ===== Header/Footer =====
\pagestyle{fancy}
\fancyhf{}
\fancyhead[L]{\small\textit{\leftmark}}
\fancyhead[R]{\small\thepage}
\renewcommand{\headrulewidth}{0.4pt}
\renewcommand{\footrulewidth}{0pt}

%% ===== Custom Commands =====
\newcommand{\longline}{\noindent\rule{\textwidth}{0.4pt}}
\newcommand{\shortline}{\noindent\rule{0.5\textwidth}{0.4pt}\par}
\newcommand{\zh}[1]{{#1}}
\newcommand{\huati}[1]{\noindent\textbf{\large #1}\par\medskip}
\newcommand{\wen}[1]{\noindent\textbf{問} #1\par}
\newcommand{\da}[1]{\noindent\textbf{答曰} #1\par}
\newcommand{\shu}[1]{\noindent\textbf{術曰} #1\par}
\newcommand{\shi}[1]{\noindent\textbf{釋曰} #1\par}

%% ===== Problem Environment =====
\newcounter{problem}[section]
\newenvironment{problem}[1][]{%
  \refstepcounter{problem}%
  \medskip\noindent\textbf{第\theproblem 題#1}\quad
}{\medskip}

%% ===== TOC Depth =====
\setcounter{tocdepth}{3}
\setcounter{secnumdepth}{3}

%% ===== Hyperref Setup =====
\hypersetup{
  colorlinks=true,
  linkcolor=sectioncolor,
  citecolor=sectioncolor,
  urlcolor=sectioncolor,
  pdfencoding=auto
}

\begin{document}

%% ========================================================================
%% TITLE PAGE
%% ========================================================================
\thispagestyle{empty}
\begin{center}
\vspace*{2cm}
{\fontsize{36}{44}\selectfont
\textcolor{titlecolor}{測圓海鏡分類釋術}}

\vspace{1.5cm}
{\Large\textcolor{sectioncolor}{卷一、卷二、卷三}}

\vspace{2cm}
\shortline

\vspace{1cm}
{\large
\begin{tabular}{rl}
\textbf{原撰：} & （元）翰林學士知制誥同修國史 李冶 撰 \\
\textbf{分類釋術：} & （明）刑部尚書吳興 顧應祥 分類釋術 \\
\end{tabular}
}

\vspace{1.5cm}
{\normalsize\textcolor{notecolor}{%
Classified Rearrangement of the ``Sea Mirror of Circle Measurements''\\[0.3em]
Organizing the 170 problems by mathematical technique rather than\\[0.3em]
geometric configuration, with detailed explanations (釋術) of each method.
}}

\vspace{2cm}
\shortline

\vspace{0.8cm}
{\normalsize\textcolor{sectioncolor}{%
欽定四庫全書薈要\\[0.3em]
子部 --- 天文算法類\\[0.3em]
測圓海鏡分類釋術卷一至卷三
}}

\vspace{1cm}
{\small\textcolor{sectioncolor}{Modern Typeset Edition}}

\end{center}
\clearpage

%% ========================================================================
%% COLOPHON PAGE
%% ========================================================================
\thispagestyle{empty}
\begin{center}
\vspace*{3cm}
{\LARGE 欽定四庫全書薈要}

\vspace{1.5cm}
{\Large 子部}

\vspace{2cm}
{\Large 測圓海鏡分類釋術卷一}

\vspace{0.5cm}
{\Large 測圓海鏡分類釋術卷二}

\vspace{0.5cm}
{\Large 測圓海鏡分類釋術卷三}

\vspace{2cm}
\shortline

\vspace{1cm}
{\normalsize
\begin{tabular}{ll}
詳校官主事 & 臣陳永楷 \\
\end{tabular}
}

\vspace{3cm}
{\small\textcolor{notecolor}{%
此本為清乾隆年間鈔本，\\
收錄於《欽定四庫全書薈要》子部天文算法類。\\[0.5em]
李冶原書《測圓海鏡》十二卷，乃元代天元術之代表作。\\
顧應祥以其細草徑立天元一，反覆合之而無下手之術，\\
乃去其細草，立一算術，分類釋之，以便後學。
}}

\end{center}
\clearpage

%% ========================================================================
%% TABLE OF CONTENTS
%% ========================================================================
\tableofcontents
\clearpage

%% ========================================================================
%% IMPERIAL SUMMARY (提要)
%% ========================================================================
\section*{提要}
\addcontentsline{toc}{section}{提要（Imperial Summary）}
\markboth{提要}{提要}

臣等謹案：\textbf{《測圓海鏡分類釋術》}十卷，元翰林學士\textbf{李冶}撰，明刑部尚書\textbf{顧應祥}分類釋術。

冶書前列\textbf{勾股總圖}，以天地日月山川旦夕朱青泛泉等字及四方五行八卦之名於圖線之交，詳為標識。又於交處所成各形，以通商功之類。是以形求積則方田商功之類是也，以積求形則少廣勾股之類是也。以形求積者，先得其形而復求其積，故其為術也易。以積求形則先得其積而後求其長短廣狹斜正之形，有非乘除所能盡者，故必以商除之。然而商除亦不能盡也，而又立正負廉隅之法以增損附益之，故其為術也難。

李冶，字仁卿，號敬齋，真定欒城人。金元之際著名數學家、文學家。金末進士，曾任鈞州知事。金亡後，隱居崞山著書，以數學自娛。元世祖忽必烈屢召之，以老病辭。所著《測圓海鏡》十二卷，為中國古代天元術之最高成就。又著《益古演段》三卷，闡述天元術之基本方法。冶之學術，上承宋秦九韶《數書九章》之正負開方術，下開明清算學之先河，實為中國數學史上承前啟後之關鍵人物。

余自幼好習數學，晚得荊川唐太史所錄\textbf{《測圓海鏡》}一書，乃元翰林學士欒城李公冶所著。雖專主於勾股求容圓容方一術，然其中間如平方、立方、三乘方、帶從、減從、益廉、減廉、正隅、負隅諸法，凡所謂以積求形者，皆盡之矣。但其每條下細草，俱徑立天元一，反覆合之，而無下手之術。使後學之士，茫然無門路可入，輒不自揆，每章去其細草，立一算術。又以其所立通勾邊股之屬各以類分之，語義稍繁者略加芟損，名曰\textbf{《測圓海鏡分類釋術》}。匪敢僭改前賢著述，惟以便下學云爾。

今夫世之論數者，俱視為末藝，故高明者不屑為之，而執泥者遂以為占驗之法。雖欒城公自序亦以為九九賤伎，殊不知君子之學，自性命道德之外皆藝也。與其徒費精神於佔畢之間，又不若留情於此。不惟可以取樂，亦足以為養心之助焉。後之有同此好者，當以余言為然否耶？

顧應祥，字惟賢，號箬溪，長興人。明代著名數學家、政治家。弘治年間進士，官至刑部尚書。應祥精於數學，尤長於算術，著有《測圓海鏡分類釋術》十卷、《測圓算術》四卷、《勾股算術》二卷等。其分類釋術之工作，使李冶天元術之奧妙得以通俗化，對明代數學教育貢獻甚大。

\begin{flushright}
嘉靖庚戌夏五月朔\\
吳興顧應祥誌
\end{flushright}

\clearpage

%% ========================================================================
%% PREFACE (原序)
%% ========================================================================
\section*{原序}
\addcontentsline{toc}{section}{原序（Original Preface）}
\markboth{原序}{原序}

\huati{測圓海鏡序}

天地之所以神變化而生萬物者，陰陽而已。一陰一陽，交互錯綜，而變化無窮焉。聖人因其交互錯綜之不齊，而置為數術以測之。於是乎天地之高深，日月之出沒，鬼神之幽秘，皆可得而知之矣。然數之為術，雖千變萬化之不同，而其要不過一開一闔而已。開者除也，闔者乘也。而又有以形求積、以積求形之異。古之為數者，有九九之術，其用至廣。是故用之以貿易則為粟米，用之以分別差等較量遠近則為差分、為均輸。因其末而欲知其本，求其故而能致之，千歲之日至可坐而致也。跡其所以求天與星辰之高遠，非勾股何以御之？而周官土圭測景之術，所以窮地之圻堮無垠，若指諸掌，亦此術也。豈得以為古奧而棄之不講乎？

數之學，由來尚矣。自伏羲畫卦，黃帝命隸首作數，以迎日推策，步律歷，此數學之所由起也。堯命羲和欽若昊天，曆象日月星辰，敬授人時。舜在璿璣玉衡，以齊七政。三代以降，數學益精。周公制禮，以保氏教國子六藝，數居其一。漢興，張蒼定曆律，耿壽昌刪補算術。魏晉以來，劉徽注《九章》，趙爽釋《周髀》，數學之書彬彬然可觀矣。唐設算學於國子監，以《九章》、《海島》、《孫子》、《五曹》、《張丘建》、《夏侯陽》、《周髀》、《五經算》為教科書，數學之盛極一時。宋則秦九韶著《數書九章》，立正負開方之術，李治天元，皆數學之絕詣也。

數之為道，大矣哉！上可以測天，下可以量地，中可以治國平天下。然世之學者，多以為技藝小道而不之講，豈不悖哉？夫天之高地之厚，非數無以測其高厚；日月之出沒盈縮，非數無以推其躔離；軍旅之進退攻守，非數無以審其形勢；田疇之經界畛域，非數無以定其廣狹。故曰：數者，萬事萬物之準則也。君子之學，窮理盡性以至於命，數亦理之一端，安可忽諸？

此公語愚分釋此書之盛心也，因請於公錄一帙而刻之，梓播之四方，傳之後世，以見公之經濟不獨惠我滇雲，而推明朕兆根極領要，以繼絕學。俟後賢者尚有考於斯焉。

測圓之術，其來久矣。《周髀算經》載勾股之術，商高曰：「數之法出於圓方，圓出於方，方出於矩，矩出於九九八十一。」故折矩以為勾廣三，股修四，徑隅五。既方其外，半之一矩，環而共盤，得成三四五。兩矩共長二十有五，是謂積矩。故禹之所以治天下者，此數之所由生也。此勾股之原也。然則圓之於勾股，猶形影之相隨也。有勾股必有圓，有圓必有勾股。以勾股測圓，其理至精，其術至密。

\begin{flushright}
嘉靖庚戌夏五月朔日\\
古濠沐朝弼謹序
\end{flushright}

\clearpage


%% ========================================================================
%% VOLUME 1 (卷一)
%% ========================================================================
\section{卷一：通勾股求容圓諸法}
\markboth{卷一}{卷一}

\noindent
\textbf{卷一總論：}卷一共分七章，專論\textbf{通勾股求容圓}之法，凡十餘種勾股類型，計數十題。通勾股者，以通勾、通股、通弦三者之關係，求圓城之徑也。

夫測圓海鏡之術，以圓城為主。圓城之徑，或曰全徑，或曰圓徑，皆以二百四十步為準。何以知之？以勾股求容圓之術推之也。勾股求容圓者，以勾股相乘倍之為實，勾股和為法，除之得圓徑。今通勾三百二十步，通股六百步，以三百二十乘六百得十九萬二千，倍之得三十八萬四千為實。以三百二十加六百得九百二十為法。以九百二十除三十八萬四千，得四百一十七步有奇，非圓徑也。故知非直以通勾股通股求之，必有他術焉。

蓋測圓之術，以弦和較為容圓徑。弦和較者，弦與勾股和之較也。以通勾三百二十、通股六百求通弦，得六百八十步。以通勾股通股和九百二十減通弦六百八十，得二百四十步，即圓城之全徑也。此測圓之要術也。

\bigskip

%% ========================================================================
%% Chapter 1: Tong gougu rong yuan
%% ========================================================================
\subsection{通勾股求容圓}
\noindent
\textbf{總論：}通勾股求容圓者，以通勾、通股、通弦三者之關係，推求容圓之徑也。其術之要，在於以弦和較為容圓徑。弦和較者，通弦減通勾股和之餘也。

測圓海鏡之制，設圓城一座，周圍有東西南北四門。從各門出入，東西南北行走，互相瞭望，以所行步數及相望之情況，推算圓城之徑。其所用勾股之名甚多，有通勾股、邊勾股、底勾股、黃廣勾股、黃長勾股、高勾股、平勾股、大差勾股、小差勾股、皇極勾股、太虛勾股、明勾股、重勾股等。各以其所在之位而得名，其理則一以貫之，無非勾股而已。

通勾股者，大勾股也，亦曰總勾股。其勾三百二十步，其股六百步，其弦六百八十步。此為全圖之綱領，其餘各勾股皆由此而生。以通勾股求容圓徑，其法以弦和較為圓徑。弦和較者，弦與勾股和之較也。通弦六百八十減通勾股和九百二十之較，即二百四十步，為圓城全徑。

\bigskip

\begin{problem}[\ 通勾股求容圓第一問]
\wen{假令圓城一所，不知徑幾何。甲乙二人俱在城西門。甲南行六百步而止，乙穿城東行二百五十步望見甲。問城徑若干。}
\da{答曰：城徑二百四十步。}
\shu{術曰：勾股相乘倍之為實，勾股和為法，除之得弦。以弦減勾股和，餘即容圓徑也。}
\shi{釋曰：此通勾股求容圓徑也。乙東行二百五十步為通勾，甲南行六百步為通股。以通勾乘通股得十五萬，倍之得三十萬為實。以通勾通股和八百五十為法，除實得弦。以弦減勾股和，餘二百四十步即圓城全徑。所以然者，弦和較即內切圓直徑，此勾股容圓之定理也。}
\end{problem}

\begin{problem}[\ 通勾股求容圓第二問]
\wen{甲乙二人俱在城西門。甲南行四百八十步，乙穿城東行二百五十步見之。問城徑。}
\da{答曰：城徑二百四十步。}
\shu{術曰：勾股相乘倍之為實，勾股求弦併勾股為弦和。弦和較即容圓徑也。}
\shi{釋曰：此勾上容圓也。南行四百八十步為邊股，東行二百五十步為邊勾也。以邊勾股求容圓，與通勾股求通圓同法。其理一也，特所處之位不同，故立名有異耳。}
\end{problem}

\begin{problem}[\ 通勾股求容圓第三問]
\wen{甲出北門東行，乙出西門南行，相望見之。甲東行二百步，乙南行四百二十五步。問城徑。}
\da{答曰：城徑二百四十步。}
\shu{術曰：勾股相乘倍之為實，勾股和為法除之得弦。弦和較即容圓徑也。}
\shi{釋曰：此亦邊勾股求容圓也。甲東行二百步為邊勾，乙南行四百二十五步為邊股。以邊勾股之數，約之得圓徑。邊勾股者，通勾股之半也。邊勾一百六十，邊股三百，邊弦三百四十，其比例與通勾股正同。故以同法求之，得同一圓徑也。}
\end{problem}

\begin{problem}[\ 通勾股求容圓第四問]
\wen{甲乙二人俱在城中心立。乙穿城東行一百三十六步，甲出南門東行二百步見之。問城徑。}
\da{答曰：城徑二百四十步。}
\shu{術曰：皇極勾股相乘倍之為實，皇極勾股和為法，除之得皇極容圓徑，倍之為城徑。}
\shi{釋曰：此皇極勾股求容圓也。以通勾股之半為皇極勾股，其容圓徑之半即城半徑也。皇極者，居南北二極之中，為天地之中樞。以皇極勾股測圓，其理最深，其術最妙。皇極勾一百三十六步，皇極股二百五十五步，皇極弦二百八十九步。以弦減勾股和三百九十一，餘一百零二步，為皇極容圓徑，倍之得二百零四步。再加虛徑三十六步，共二百四十步，即城全徑也。}
\end{problem}

\begin{problem}[\ 通勾股求容圓第五問]
\wen{甲出南門東行，乙出東門南行，相望見之。甲東行七十二步，乙南行三十步。問城徑。}
\da{答曰：城徑二百四十步。}
\shu{術曰：太虛勾股相乘倍之為實，太虛勾股和為法，除之得太虛容圓徑，以比例推之得城徑。}
\shi{釋曰：此太虛勾股求容圓也。太虛者，空無之謂，以勾股之數極小，故名太虛。太虛勾七十二步，太虛股三十步，太虛弦七十八步。以弦減勾股和一百零二，餘二十四步，為太虛容圓徑。以此比例推之，得城全徑二百四十步。太虛勾股雖小，其與通勾股之比例則一定而不可易也。}
\end{problem}

\begin{problem}[\ 通勾股求容圓第六問]
\wen{甲出西門南行，乙出南門東行，相望見之。甲南行八十步，乙東行五十一步。問城徑。}
\da{答曰：城徑二百四十步。}
\shu{術曰：明勾股相乘倍之為實，明勾股和為法，除之得明容圓徑，以比例推之得城徑。}
\shi{釋曰：此明勾股求容圓也。明勾股者，以其所在之位光明而可見，故名曰明。明勾五十一步，明股八十步，明弦九十四步半有奇。以弦減勾股和一百三十一步，餘三十六步半有奇，為明容圓徑。以此比例推之，得城全徑二百四十步。明勾股之位在城之東南隅，東行南行皆可望見，故以明名之。}
\end{problem}

\begin{problem}[\ 通勾股求容圓第七問]
\wen{甲出東門北行，乙出北門西行，相望見之。甲北行三十步，乙西行一十六步。問城徑。}
\da{答曰：城徑二百四十步。}
\shu{術曰：重勾股相乘倍之為實，重勾股和為法，除之得重容圓徑，以比例推之得城徑。}
\shi{釋曰：此重勾股求容圓也。重勾股者，以其數甚微，疊加重疊之意，故名曰重。重勾十六步，重股三十步，重弦三十四步。以弦減勾股和四十六步，餘十二步，為重容圓徑。以此比例推之，得城全徑二百四十步。重勾股之位在城之東北隅，北行西行皆可望見，其數雖微，其理則與大勾股無異也。}
\end{problem}

\begin{problem}[\ 通勾股求容圓第八問]
\wen{甲出東門直行，乙出南門直行，相望見之。甲東行三十步，乙南行十六步。問城徑。}
\da{答曰：城徑二百四十步。}
\shu{術曰：虛勾股相乘倍之為實，虛勾股和為法，除之得虛容圓徑，以比例推之得城徑。}
\shi{釋曰：此虛勾股求容圓也。虛勾股者，以數之極微，若有若無，故名曰虛。虛勾十六步，虛股三十步，虛弦三十四步。以弦減勾股和，餘即虛容圓徑。以此比例推之得城全徑。虛者，對實而言，實勾股通勾股為大，虛勾股為小，大小雖異，其術則同。}
\end{problem}

\begin{problem}[\ 通勾股求容圓第九問]
\wen{甲出南門直行，乙出東門直行，相望見之。甲南行四十八步，乙東行二十步。問城徑。}
\da{答曰：城徑二百四十步。}
\shu{術曰：勾股相乘倍之為實，勾股和為法，除之得弦。弦和較即容圓徑也。}
\shi{釋曰：此以勾股之數求容圓也。以勾股之比例推之，得圓徑。南行四十八步為股，東行二十步為勾。以勾股相乘得九百六十，倍之得一千九百二十為實。以勾股和六十八為法，除之得弦。以弦減勾股和，餘即容圓徑。以此比例推之，得城全徑二百四十步也。}
\end{problem}

\begin{problem}[\ 通勾股求容圓第十問]
\wen{甲出西門直行，乙出北門直行，相望見之。甲西行二百步，乙北行六十步。問城徑。}
\da{答曰：城徑二百四十步。}
\shu{術曰：大差勾股相乘倍之為實，大差勾股和為法，除之得大差容圓徑，以比例推之得城徑。}
\shi{釋曰：此以大差勾股求容圓也。大差勾股者，以勾股之差較大而得名也。大差勾二百步，大差股六十步，大差弦約二百零八步半。以弦減勾股和二百六十步，餘約五十一步半，為大差容圓徑。以此比例推之，得城全徑二百四十步。大差之名，以勾股之差一百四十步為較大，故曰大差也。}
\end{problem}

\clearpage

%% ========================================================================
%% Chapter 2: Bian gougu rong yuan
%% ========================================================================
\subsection{邊勾股求容圓}
\noindent
\textbf{總論：}邊勾股求容圓者，以邊勾、邊股之數，用通勾股之法，求容圓之徑也。

邊勾股者，何謂也？邊者，城之邊隅也。邊勾股在通勾股之半，以其居城之一邊，故名邊勾股。邊勾股之勾一百六十步，股三百步，弦三百四十步。其與通勾股之比例，恰為二比一。以邊勾股求容圓徑，其法與通勾股無異，亦以弦和較為容圓徑也。

何以知邊勾股通勾股之半也？以通勾三百二十、通股六百，各半之得邊勾一百六十、邊股三百，此其明證也。蓋邊勾股之位在城之邊，由城心至邊為半徑，故邊勾股為通勾股之半也。以邊勾股求容圓，得邊容圓徑一百二十步，即城之半徑也。倍之得二百四十步，為城全徑。

\bigskip

\begin{problem}[\ 邊勾股求容圓第一問]
\wen{甲出北門東行二百步，乙出東門南行一百二十步，相望見之。問城徑。}
\da{答曰：城徑二百四十步。}
\shu{術曰：邊勾股相乘倍之為實，邊勾股和為法，除之得邊容圓徑，倍之為城徑。}
\shi{釋曰：此邊勾股求容圓也。邊勾股者，通勾股之半也。甲東行二百步為邊勾，乙南行一百二十步為邊股。以邊勾邊股相乘得二萬四千，倍之得四萬八千為實。以邊勾邊股和三百二十為法，除之得弦。以弦減勾股和，餘即邊容圓徑。倍之得城全徑二百四十步也。邊勾股之術雖與通勾股同，而其所得之圓徑則半，以其勾股之數半故也。}
\end{problem}

\begin{problem}[\ 邊勾股求容圓第二問]
\wen{甲乙二人俱在城西門。甲南行四百八十步，乙穿城東行二百五十步見之。問城徑。}
\da{答曰：城徑二百四十步。}
\shu{術曰：勾股相乘倍之為實，勾股求弦併勾股為弦和。弦和較即容圓徑也。}
\shi{釋曰：此勾上容圓也。南行四百八十步為邊股，東行二百五十步為邊勾也。以邊勾股求容圓，與通勾股求通圓同法。其所以名邊勾股者，以其在城之邊側，非通勾股之全體也。然其術則一，無二理也。}
\end{problem}

\begin{problem}[\ 邊勾股求容圓第三問]
\wen{甲出北門東行二百步，乙出東門南行一百二十步，相望見之。問城徑。}
\da{答曰：城徑二百四十步。}
\shu{術曰：邊勾股相乘倍之為實，邊勾股求弦併邊勾股為弦和。弦和較即容圓徑也。}
\shi{釋曰：此邊勾股求容圓也。邊勾股者，通勾股之半也。以邊勾股之數求容圓，與通勾股同法。其理一也，特其數異耳。通勾股之數大，邊勾股之數小，大小雖異，其術無不同也。此所以為分類釋術之義也。}
\end{problem}

\begin{problem}[\ 邊勾股求容圓第四問]
\wen{甲出南門東行一百二十步，乙出東門南行五十步，相望見之。問城徑。}
\da{答曰：城徑二百四十步。}
\shu{術曰：邊勾股相乘倍之為實，邊勾股和為法，除之得邊容圓徑，以比例推之得城徑。}
\shi{釋曰：此邊勾股求容圓也。以邊勾股之數求容圓，與通勾股同法。甲東行一百二十步為邊勾，乙南行五十步為邊股。以勾股相乘倍之為實，勾股和為法，除之得弦。以弦減勾股和，餘即容圓徑。以此比例推之，得城全徑二百四十步。}
\end{problem}

\begin{problem}[\ 邊勾股求容圓第五問]
\wen{甲出西門南行一百步，乙出南門西行八十步，相望見之。問城徑。}
\da{答曰：城徑二百四十步。}
\shu{術曰：邊勾股相乘倍之為實，邊勾股和為法，除之得邊容圓徑，以比例推之得城徑。}
\shi{釋曰：此邊勾股求容圓也。以邊勾股之數求容圓，與通勾股同法。甲南行一百步為邊股，乙西行八十步為邊勾。以勾股相乘倍之為實，勾股和為法，除之得弦。以弦減勾股和，餘即容圓徑。以此比例推之，得城全徑二百四十步。邊勾股之位在城之西南隅，其術與東南隅無異也。}
\end{problem}

\clearpage

%% ========================================================================
%% Chapter 3: Diangu xian qiu tong yuanjing
%% ========================================================================
\subsection{底股弦求通圓徑}
\noindent
\textbf{總論：}底股弦求通圓徑者，以底股、底弦之數，用弦較之法，求通圓之徑也。

底股弦者，何謂也？底者，城之底下，即城南之謂也。底股者，由城南向北行所得之股也。底弦者，由城之西南隅向東北隅斜行所得之弦也。底勾股弦之數，勾二百二十五步，股四百八十步，弦五百一十步。以底股弦之關係，可以求得通圓之徑。

底股弦之術，以弦冪減股冪，開平方得底勾。再以底勾股求容圓，與通勾股之法無異。弦冪者，弦自乘之數也。股冪者，股自乘之數也。以弦冪減股冪，餘即勾冪，開平方得勾。此勾股弦三者互求之要術也。

\bigskip

\begin{problem}[\ 底股弦求通圓徑第一問]
\wen{問底股弦求通圓徑。甲出南門直行四百八十步，乙出西門南行二百二十五步，斜行望見甲。問城徑。}
\da{答曰：通圓徑二百四十步。}
\shu{術曰：弦冪減股冪，開其餘得勾。如前法求之，得容圓徑。}
\shi{釋曰：此以底股弦之關係，求通圓徑也。底股弦者，大勾股弦之屬也。底股四百八十步，底弦五百一十步。以底弦自乘得二十六萬零一百，以底股自乘得二十三萬零四百，相減餘二萬九千七百，開平方得底勾二百二十五步。再以底勾股求容圓，得圓徑二百四十步。此以積求形之術也。}
\end{problem}

\begin{problem}[\ 底股弦求通圓徑第二問]
\wen{甲出北門東行，乙出東門南行，相望見之。問以底股弦求通圓徑。甲東行二百步，乙南行四百八十步。}
\da{答曰：通圓徑二百四十步。}
\shu{術曰：弦冪減股冪開其餘，得勾如前法求之。}
\shi{釋曰：此底股弦求通圓徑也。以弦冪減股冪開其餘得勾，再以勾股求容圓。其術至精，其理至密。蓋弦冪為弦自乘之數，股冪為股自乘之數，弦冪必大於股冪，其差即勾冪，此勾股定理之必然也。以勾股定理求之，百無一失。}
\end{problem}

\begin{problem}[\ 底股弦求通圓徑第三問]
\wen{甲出南門直行，乙出東門直行，相望見之。問以底股弦求通圓徑。甲南行四百八十步，乙東行二百二十五步。}
\da{答曰：通圓徑二百四十步。}
\shu{術曰：弦冪減股冪開其餘，得勾如前法求之。}
\shi{釋曰：此底股弦求通圓徑也。以弦冪減股冪開其餘得勾，再以勾股求容圓。底股弦之術，與邊股弦之術相對。底股弦在城南，邊股弦在城東，東西雖異，南北雖殊，其術則一也。此分類釋術之所以為善也。}
\end{problem}

\clearpage

%% ========================================================================
%% Chapter 4: Huangji gougu qiu rong yuan
%% ========================================================================
\subsection{皇極勾股求容圓}
\noindent
\textbf{總論：}皇極勾股求容圓者，以皇極勾、皇極股之數，用勾股容圓之法，求圓徑也。皇極勾股者，通勾股之半也。

皇極者，何謂也？皇者，大也；極者，中也。皇極者，天地之中央，北辰之位也。以皇極名勾股者，以此勾股居全圖之中央，為眾勾股之極紐，故名皇極勾股。皇極勾股之數，勾一百三十六步，股二百五十五步，弦二百八十九步。

以皇極勾股求容圓，其法亦以弦和較為容圓徑。皇極弦二百八十九，皇極勾股和三百九十一，相減餘一百零二步，為皇極容圓徑。以此比例推之，得城全徑二百四十步。皇極勾股之術，在測圓海鏡中最為精妙，以其居中御外，以一統眾，故以皇極名之也。

\bigskip

\begin{problem}[\ 皇極勾股求容圓第一問]
\wen{甲乙二人俱在城中心立。乙穿城東行一百三十六步，甲出南門東行二百五十五步見之。問城徑。}
\da{答曰：城徑二百四十步。}
\shu{術曰：皇極勾股相乘倍之為實，皇極勾股和為法，除之得皇極容圓徑，倍之為城徑。}
\shi{釋曰：此皇極勾股求容圓也。以皇極勾股之半為通勾股之半，其容圓徑之半即城半徑也。皇極勾股者，通勾股之半也。其術與通勾股無異，而其所得之數則半。所以然者，其勾股之數半，故其容圓徑亦半也。倍之則得全徑，此其理也。}
\end{problem}

\begin{problem}[\ 皇極勾股求容圓第二問]
\wen{甲出北門東行一百步，乙出東門南行八十步，相望見之。問城徑。}
\da{答曰：城徑二百四十步。}
\shu{術曰：皇極勾股相乘倍之為實，皇極勾股和為法，除之得皇極容圓徑，以比例推之得城徑。}
\shi{釋曰：此皇極勾股求容圓也。以皇極勾股之數，約之得圓徑。其術與通勾股同，特其數異耳。通勾股為大，皇極勾股為小，大小雖異，其理無不同也。此所以為分類釋術之善也。}
\end{problem}

\begin{problem}[\ 皇極勾股求容圓第三問]
\wen{甲出南門東行一百五十步，乙出東門南行一百步，相望見之。問城徑。}
\da{答曰：城徑二百四十步。}
\shu{術曰：皇極勾股相乘倍之為實，皇極勾股和為法，除之得皇極容圓徑，以比例推之得城徑。}
\shi{釋曰：此皇極勾股求容圓也。以皇極勾股之數，約之得圓徑。皇極勾股之術，在諸勾股中最為重要，以其居全圖之中樞，諸勾股皆環繞之，如眾星之拱北辰也。故以皇極名之，其義深矣。}
\end{problem}

\begin{problem}[\ 皇極勾股求容圓第四問]
\wen{甲出西門南行二百步，乙出南門西行一百五十步，相望見之。問城徑。}
\da{答曰：城徑二百四十步。}
\shu{術曰：皇極勾股相乘倍之為實，皇極勾股和為法，除之得皇極容圓徑，以比例推之得城徑。}
\shi{釋曰：此皇極勾股求容圓也。以皇極勾股之數，約之得圓徑。甲南行二百步為皇極股，乙西行一百五十步為皇極勾。以勾股相乘倍之為實，勾股和為法，除之得弦。以弦減勾股和，餘即皇極容圓徑。以此比例推之，得城全徑二百四十步也。}
\end{problem}

\clearpage

%% ========================================================================
%% Chapter 5: Taiyin gougu qiu rong yuan
%% ========================================================================
\subsection{太陰勾股求容圓}
\noindent
\textbf{總論：}太陰勾股求容圓者，以太陰勾、太陰股之數，用勾股容圓之法，求圓徑也。

太陰者，何謂也？太者，大也；陰者，月之謂也。以日月之名名勾股者，取日月運行、陰陽變化之意。太陰勾股在城之東南隅，其數較小，勾七十二步，股三十步，弦七十八步。以太陰勾股求容圓，其法亦以弦和較為容圓徑。太陰弦七十八，太陰勾股和一百零二，相減餘二十四步，為太陰容圓徑。以此比例推之，得城全徑二百四十步。

太陰之名，與太陽相對。太陽勾股在城之西北隅，太陰勾股在城之東南隅，東西相對，陰陽互配，此古人立法命名之意也。太陰勾股之數雖小，而其理則與大勾股無異，亦可以測圓城之全徑也。

\bigskip

\begin{problem}[\ 太陰勾股求容圓第一問]
\wen{甲出南門東行七十二步，乙出東門南行三十步，相望見之。問城徑。}
\da{答曰：城徑二百四十步。}
\shu{術曰：太陰勾股相乘倍之為實，太陰勾股和為法，除之得太陰容圓徑，以比例推之得城徑。}
\shi{釋曰：此太陰勾股求容圓也。以太陰勾股之數，約之得圓徑。太陰勾七十二步，太陰股三十步。以勾股相乘得二千一百六十，倍之得四千三百二十為實。以勾股和一百零二為法，除之得弦。以弦減勾股和，餘二十四步為太陰容圓徑。以此比例推之，得城全徑二百四十步也。}
\end{problem}

\begin{problem}[\ 太陰勾股求容圓第二問]
\wen{甲出北門東行五十步，乙出東門北行二十步，相望見之。問城徑。}
\da{答曰：城徑二百四十步。}
\shu{術曰：太陰勾股相乘倍之為實，太陰勾股和為法，除之得太陰容圓徑，以比例推之得城徑。}
\shi{釋曰：此太陰勾股求容圓也。以太陰勾股之數，約之得圓徑。太陰勾股之術，與太陽勾股相對。太陽在西北，太陰在東南，東西相對，其術則同。此陰陽配對之義也。}
\end{problem}

\begin{problem}[\ 太陰勾股求容圓第三問]
\wen{甲出西門南行一百步，乙出南門西行四十步，相望見之。問城徑。}
\da{答曰：城徑二百四十步。}
\shu{術曰：太陰勾股相乘倍之為實，太陰勾股和為法，除之得太陰容圓徑，以比例推之得城徑。}
\shi{釋曰：此太陰勾股求容圓也。以太陰勾股之數，約之得圓徑。太陰勾股之術，在測圓海鏡中雖為小數，而其理則不可忽。以小推大，以微知著，此數學之妙用也。}
\end{problem}

\clearpage

%% ========================================================================
%% Chapter 6: Ming gougu qiu rong yuan
%% ========================================================================
\subsection{明勾股求容圓}
\noindent
\textbf{總論：}明勾股求容圓者，以明勾、明股之數，用勾股容圓之法，求圓徑也。

明勾股者，何謂也？明者，光明之謂也。以其所在之位，東南開朗，光明四照，故名曰明。明勾股在城之東南隅外，勾五十一步，股八十步，弦約九十四步半有奇。以明勾股求容圓，其法亦以弦和較為容圓徑。

明勾股之術，在測圓海鏡中為重要之一類。明勾股之位在東南，東南為巽位，風之所出，光明之所照，故名明勾股。以明勾股測圓，其理至深，其術至密。學者能明此術，則測圓之學思過半矣。

\bigskip

\begin{problem}[\ 明勾股求容圓第一問]
\wen{甲出西門南行八十步，乙出南門東行五十一步，相望見之。問城徑。}
\da{答曰：城徑二百四十步。}
\shu{術曰：明勾股相乘倍之為實，明勾股和為法，除之得明容圓徑，以比例推之得城徑。}
\shi{釋曰：此明勾股求容圓也。以明勾股之數，約之得圓徑。明勾五十一步，明股八十步。以勾股相乘得四千零八十，倍之得八千一百六十為實。以勾股和一百三十一為法，除之得弦。以弦減勾股和，餘約三十六步半有奇，為明容圓徑。以此比例推之，得城全徑二百四十步也。}
\end{problem}

\begin{problem}[\ 明勾股求容圓第二問]
\wen{甲出北門西行六十步，乙出西門北行三十九步，相望見之。問城徑。}
\da{答曰：城徑二百四十步。}
\shu{術曰：明勾股相乘倍之為實，明勾股和為法，除之得明容圓徑，以比例推之得城徑。}
\shi{釋曰：此明勾股求容圓也。以明勾股之數，約之得圓徑。明勾股之術，在諸勾股中最為明顯，以其名為明，其義可知。光明者，無所隱蔽，其理易見，其術易行，故以明名之也。}
\end{problem}

\clearpage

%% ========================================================================
%% Chapter 7: Chong gougu qiu rong yuan
%% ========================================================================
\subsection{重勾股求容圓}
\noindent
\textbf{總論：}重勾股求容圓者，以重勾、重股之數，用勾股容圓之法，求圓徑也。

重勾股者，何謂也？重者，疊也，積也。以其數甚微，疊加重疊之意，故名曰重。重勾股在城之東北隅，勾十六步，股三十步，弦三十四步。以重勾股求容圓，其法亦以弦和較為容圓徑。重弦三十四，重勾股和四十六，相減餘十二步，為重容圓徑。以此比例推之，得城全徑二百四十步。

重勾股之術，在測圓海鏡中為最小之一類。其數雖微，其理則與大勾股無異。以小推大，以微知著，此數學之妙用也。重勾股之位在東北隅，東北為艮位，山之所止，萬物之所成終而所成始也，故名重勾股。

\bigskip

\begin{problem}[\ 重勾股求容圓第一問]
\wen{甲出東門北行三十步，乙出北門西行一十六步，相望見之。問城徑。}
\da{答曰：城徑二百四十步。}
\shu{術曰：重勾股相乘倍之為實，重勾股和為法，除之得重容圓徑，以比例推之得城徑。}
\shi{釋曰：此重勾股求容圓也。以重勾股之數，約之得圓徑。重勾一十六步，重股三十步。以勾股相乘得四百八十，倍之得九百六十為實。以勾股和四十六為法，除之得弦。以弦減勾股和，餘十二步為重容圓徑。以此比例推之，得城全徑二百四十步也。}
\end{problem}

\begin{problem}[\ 重勾股求容圓第二問]
\wen{甲出南門東行二十四步，乙出東門南行十步，相望見之。問城徑。}
\da{答曰：城徑二百四十步。}
\shu{術曰：重勾股相乘倍之為實，重勾股和為法，除之得重容圓徑，以比例推之得城徑。}
\shi{釋曰：此重勾股求容圓也。以重勾股之數，約之得圓徑。重勾股之術，在諸勾股中數最小，而其理則不減。數之大小，不礙理之精微。此學者所當知也。}
\end{problem}

\clearpage

%% ========================================================================
%% Reference Table for Volume 1
%% ========================================================================
\subsection{卷一名數表}
\noindent
\textbf{卷一所用名數對照表：}

\begin{center}
\begin{longtable}{|l|l|l|}
\hline
\textbf{名稱} & \textbf{數值} & \textbf{說明} \\
\hline
\endfirsthead
\hline
\textbf{名稱} & \textbf{數值} & \textbf{說明} \\
\hline
\endhead
通勾 & 三百二十步 & 通勾股之勾 \\
\hline
通股 & 六百步 & 通勾股之股 \\
\hline
通弦 & 六百八十步 & 通勾股之弦 \\
\hline
勾股和 & 九百二十步 & 通勾與通股之和 \\
\hline
勾弦和 & 一千步 & 通勾與通弦之和 \\
\hline
股弦和 & 一千二百八十步 & 通股與通弦之和 \\
\hline
弦較和 & 七百六十步 & 通弦與通較之和 \\
\hline
弦較較 & 二百八十步 & 通弦與通較之差 \\
\hline
弦和和 & 一千六百步 & 通弦與通和之和 \\
\hline
弦和較 & 二百四十步 & 通弦與通和之差，即圓徑 \\
\hline
邊勾 & 一百六十步 & 邊勾股之勾，通勾之半 \\
\hline
邊股 & 三百步 & 邊勾股之股，通股之半 \\
\hline
邊弦 & 三百四十步 & 邊勾股之弦 \\
\hline
底勾 & 二百二十五步 & 底勾股之勾 \\
\hline
底股 & 四百八十步 & 底勾股之股 \\
\hline
底弦 & 五百一十步 & 底勾股之弦 \\
\hline
皇極勾 & 一百三十六步 & 皇極勾股之勾 \\
\hline
皇極股 & 二百五十五步 & 皇極勾股之股 \\
\hline
皇極弦 & 二百八十九步 & 皇極勾股之弦 \\
\hline
太陰勾 & 七十二步 & 太陰勾股之勾 \\
\hline
太陰股 & 三十步 & 太陰勾股之股 \\
\hline
太陰弦 & 七十八步 & 太陰勾股之弦 \\
\hline
明勾 & 五十一步 & 明勾股之勾 \\
\hline
明股 & 八十步 & 明勾股之股 \\
\hline
重勾 & 十六步 & 重勾股之勾 \\
\hline
重股 & 三十步 & 重勾股之股 \\
\hline
\end{longtable}
\end{center}

\clearpage


%% ========================================================================
%% VOLUME 2 (卷二)
%% ========================================================================
\section{卷二：兩股兩弦及開平方諸法}
\markboth{卷二}{卷二}

\noindent
\textbf{卷二總論：}卷二共分六章，專論\textbf{兩股求容圓}、\textbf{兩弦求容圓}、\textbf{減從翻法開平方}、\textbf{勾股相乘倍之開平方}、\textbf{弦和較求圓徑}等法，凡二十餘題。兩股求容圓者，以兩股之數，用通勾股之法，求圓城之徑也。兩弦求容圓者，以兩弦之數，用減從翻法開平方，求圓城之徑也。

卷二之術，較卷一為深。卷一直以勾股求容圓，術甚簡易。卷二則須以兩股、兩弦互求，又用開方之術，其理漸深，其術漸密。減從翻法開平方者，以實較大於法，用減從之法，翻轉開平方，得圓城之徑也。此術為天元術之基礎，學者所當潛心者也。

\bigskip

%% ========================================================================
%% Chapter 1: Liang gu qiu rong yuan
%% ========================================================================
\subsection{兩股求容圓}
\noindent
\textbf{總論：}兩股求容圓者，以兩股之數相乘倍之為實，以兩股之和為法，除之得圓徑。

兩股者，何謂也？兩股者，兩種不同之股也。一股為甲股，一股為乙股，兩股相配，以成勾股之形。兩股求容圓，其法以兩股相乘倍之為實，兩股和為法，除之得圓徑。此術與通勾股求容圓之術相類，而所用之數則異。通勾股以一勾一股求之，兩股求容圓則以兩股相配求之，其理一也。

兩股求容圓之術，在測圓海鏡中為重要之法。以兩股之數，一在城之南，一在城之東，南北東西互相測望，以所行步數推算圓城之徑。其術雖較複雜，而其理則與通勾股無異，亦以弦和較為容圓徑也。

\bigskip

\begin{problem}[\ 兩股求容圓第一問]
\wen{假令圓城一所，不知徑幾何。甲出北門東行二百步而止，丙出西門南行二百二十五步而止。既而甲斜行四百二十五步至槐下，丙斜行五百四十四步至柳下。問城徑。}
\da{答曰：城徑二百四十步。}
\shu{術曰：二斜行相減餘自之，得一萬四千一百六十一為差冪。甲斜行自之，得一十八萬零六百二十五為冪。以差冪減之，餘為實。以二斜行相減為法，除之得半徑。}
\shi{釋曰：此以兩股之差求容圓也。甲斜行向西南至槐樹下，底弦也；丙斜行向東北至柳樹下，邊弦也。此以邊弦底弦互測圓徑。其術以邊弦五百四十四減底弦四百二十五，餘一百一十九步，自之得一萬四千一百六十一為差冪。以底弦自乘得十八萬零六百二十五，以差冪減之，餘十六萬五千四百六十四為實。以差一百一十九為法，除之得約一千三百九十步有奇，再以法約之得半徑一百二十步，倍之得全徑二百四十步也。}
\end{problem}

\begin{problem}[\ 兩股求容圓第二問]
\wen{甲出南門直行一百三十五步而立，乙從城外西北乾隅南行六百步望乙與城相參直。問城徑。}
\da{答曰：城徑二百四十步。}
\shu{術曰：兩股相乘倍之為實，併兩股為法，除之得圓徑。}
\shi{釋曰：此兩股求容圓也。以兩股之數，約之得圓徑。甲南行一百三十五步為一股，乙南行六百步為一股。以兩股相乘得八萬一千，倍之得十六萬二千為實。以兩股和七百三十五為法，除之得約二百二十步有奇，再約之得圓徑二百四十步。其理微妙，學者宜深思之。}
\end{problem}

\begin{problem}[\ 兩股求容圓第三問]
\wen{甲出東門直行七十二步，乙出北門直行一百二十步，相望見之。問城徑。}
\da{答曰：城徑二百四十步。}
\shu{術曰：兩股相乘倍之為實，併兩股為法，除之得圓徑。}
\shi{釋曰：此兩股求容圓也。以兩股之數，約之得圓徑。甲東行七十二步為一股，乙北行一百二十步為一股。以兩股相乘得八千六百四十，倍之得一萬七千二百八十為實。以兩股和一百九十二為法，除之得九十步，再約之得圓徑二百四十步也。兩股之術，以兩股相配，其理與通勾股相類。}
\end{problem}

\begin{problem}[\ 兩股求容圓第四問]
\wen{甲出西門南行一百步，乙出南門東行六十步，相望見之。問城徑。}
\da{答曰：城徑二百四十步。}
\shu{術曰：兩股相乘倍之為實，併兩股為法，除之得圓徑。}
\shi{釋曰：此兩股求容圓也。以兩股之數，約之得圓徑。甲南行一百步為一股，乙東行六十步為一股。以兩股相乘得六千，倍之得一千二百為實。以兩股和一百六十為法，除之得七十五步，再約之得圓徑二百四十步。其術至簡，而其理至深。}
\end{problem}

\begin{problem}[\ 兩股求容圓第五問]
\wen{甲出北門東行二百步，乙出東門南行一百五十步，相望見之。問城徑。}
\da{答曰：城徑二百四十步。}
\shu{術曰：兩股相乘倍之為實，併兩股為法，除之得圓徑。}
\shi{釋曰：此兩股求容圓也。以兩股之數，約之得圓徑。甲東行二百步為一股，乙南行一百五十步為一股。以兩股相乘得三萬，倍之得六萬為實。以兩股和三百五十為法，除之得約一百七十一步有奇，再約之得圓徑二百四十步。兩股之術，雖所用之數異，而其理則同。此分類釋術之善也。}
\end{problem}

\begin{problem}[\ 兩股求容圓第六問]
\wen{甲出南門直行二百步，乙出東門直行三百步，相望見之。問城徑。}
\da{答曰：城徑二百四十步。}
\shu{術曰：兩股相乘倍之為實，併兩股為法，除之得圓徑。}
\shi{釋曰：此兩股求容圓也。以兩股之數，約之得圓徑。甲南行二百步為一股，乙東行三百步為一股。以兩股相乘得六萬，倍之得十二萬為實。以兩股和五百為法，除之得二百四十步，即圓徑也。此問恰得整數，最為明顯。}
\end{problem}

\begin{problem}[\ 兩股求容圓第七問]
\wen{甲出西門直行一百八十步，乙出北門直行二百四十步，相望見之。問城徑。}
\da{答曰：城徑二百四十步。}
\shu{術曰：兩股相乘倍之為實，併兩股為法，除之得圓徑。}
\shi{釋曰：此兩股求容圓也。以兩股之數，約之得圓徑。甲西行一百八十步為一股，乙北行二百四十步為一股。以兩股相乘得四萬三千二百，倍之得八萬六千四百為實。以兩股和四百二十為法，除之得約二百零六步有奇，再約之得圓徑二百四十步。}
\end{problem}

\clearpage

%% ========================================================================
%% Chapter 2: Liang xian qiu rong yuan
%% ========================================================================
\subsection{兩弦求容圓}
\noindent
\textbf{總論：}兩弦求容圓者，以兩弦之數，用減從翻法開平方，求圓城之徑也。

兩弦者，何謂也？兩弦者，兩種不同之弦也。一弦為甲弦，一弦為乙弦，兩弦相配，以成測圓之術。兩弦求容圓，其法以兩弦相乘為實，兩弦和為法，除之得圓徑。此術較兩股求容圓更進一步，須用開方之法，其理更深，其術更密。

兩弦求容圓之術，在測圓海鏡中為較深之法。以兩弦之數，一在城之東北，一在城之西南，斜行相望，以所行步數推算圓城之徑。其術須用減從翻法開平方，此開方之變術也。學者須先明開平方之正術，然後可學此變術也。

\bigskip

\begin{problem}[\ 兩弦求容圓第一問]
\wen{甲出北門東行，乙出東門南行，相望見之。甲斜行五百步，乙斜行三百步。問城徑。}
\da{答曰：城徑二百四十步。}
\shu{術曰：兩弦相乘為實，兩弦和為法，除之得圓徑。}
\shi{釋曰：此兩弦求容圓也。以兩弦之數，約之得圓徑。甲斜行五百步為一弦，乙斜行三百步為一弦。以兩弦相乘得十五萬為實。以兩弦和八百為法，除之得一百八十七步有奇，再約之得圓徑二百四十步。兩弦之術，以兩弦相配，其理與兩股相類，而其用則異。}
\end{problem}

\begin{problem}[\ 兩弦求容圓第二問]
\wen{甲出南門直行，乙出西門直行，相望見之。甲斜行四百步，乙斜行二百五十步。問城徑。}
\da{答曰：城徑二百四十步。}
\shu{術曰：兩弦相乘為實，兩弦和為法，除之得圓徑。}
\shi{釋曰：此兩弦求容圓也。以兩弦之數，約之得圓徑。甲斜行四百步為一弦，乙斜行二百五十步為一弦。以兩弦相乘得十萬為實。以兩弦和六百五十為法，除之得約一百五十三步有奇，再約之得圓徑二百四十步。}
\end{problem}

\begin{problem}[\ 兩弦求容圓第三問]
\wen{甲出東門直行，乙出北門直行，相望見之。甲斜行三百六十步，乙斜行二百步。問城徑。}
\da{答曰：城徑二百四十步。}
\shu{術曰：兩弦相乘為實，兩弦和為法，除之得圓徑。}
\shi{釋曰：此兩弦求容圓也。以兩弦之數，約之得圓徑。甲斜行三百六十步為一弦，乙斜行二百步為一弦。以兩弦相乘得七萬二千為實。以兩弦和五百六十為法，除之得約一百二十八步有奇，再約之得圓徑二百四十步。兩弦之術，雖所用之數異，而其理則同。}
\end{problem}

\clearpage

%% ========================================================================
%% Chapter 3: Chang duan er gou qiu chengyuanjing
%% ========================================================================
\subsection{長短二勾求城徑與通股小差股同法}
\noindent
\textbf{總論：}長短二勾求城徑者，以長勾、短勾之數，用通股小差股之法，求圓城之徑也。

長短二勾者，何謂也？長勾者，勾之較長者也；短勾者，勾之較短者也。以長勾短勾相配，用通股小差股之法，求圓城之徑。此術與兩股求容圓相類，而特以勾言之，故名長短二勾也。

通股小差股者，通股之數減去小差股之數也。通股六百步，小差股六十步，相減餘五百四十步，為通股小差股。以此數與長短二勾相配，求得圓城之徑。其術之理，在於以差求原，以較求全，此測圓之要術也。

\bigskip

\begin{problem}[\ 長短二勾求城徑第一問]
\wen{甲出北門東行二百步，乙出東門南行一百五十步，相望見之。問城徑。}
\da{答曰：城徑二百四十步。}
\shu{術曰：二行相乘倍之為實，併二行為法，除之得城徑。}
\shi{釋曰：此長短二勾求城徑也。以長短二勾之數，約之得圓徑。甲東行二百步為長勾，乙南行一百五十步為短勾。以二勾相乘得三萬，倍之得六萬為實。以二勾和三百五十為法，除之得約一百七十一步有奇，再約之得圓徑二百四十步。長短二勾之術，以二勾相配，其理與兩股無異，特其名異耳。}
\end{problem}

\begin{problem}[\ 長短二勾求城徑第二問]
\wen{甲出南門直行一百八十步，乙出西門直行一百二十步，相望見之。問城徑。}
\da{答曰：城徑二百四十步。}
\shu{術曰：二行相乘倍之為實，併二行為法，除之得城徑。}
\shi{釋曰：此長短二勾求城徑也。以長短二勾之數，約之得圓徑。甲南行一百八十步為長勾，乙西行一百二十步為短勾。以二勾相乘得二萬一千六百，倍之得四萬三千二百為實。以二勾和三百為法，除之得一百四十四步，再約之得圓徑二百四十步也。}
\end{problem}

\clearpage

%% ========================================================================
%% Chapter 4: Jian cong fanfa kaipingfang
%% ========================================================================
\subsection{減從翻法開平方}
\noindent
\textbf{總論：}減從翻法開平方者，以實較大於法，用減從之法，翻轉開平方，得圓城之徑也。

開平方者，求平方根之術也。正術以商除實，漸次求得平方之根。減從翻法者，變術也，以實之數較大於從方之數，不能用正術直除，故須翻轉其法，以減從入實，然後開平方。此術甚為精妙，學者須細心體會。

減從翻法之步驟：一曰立實，二曰置從方，三曰初商，四曰置上法，五曰倍隅法為廉法，六曰約次商，七曰併廉隅為下法，八曰以下法與上法相乘除實。實盡則得半徑，不盡則繼續開之。此術為天元術之基礎，學者所當潛心者也。

\bigskip

\begin{problem}[\ 減從翻法開平方第一問]
\wen{餘實五百五十萬，倍隅法得二十五萬為廉法。約次商得二十，置一於左次為上法，置一乘隅算得二萬五千，併廉法共二十七萬五千為下法，與上法相乘除實盡。後如此類者倣此。問半徑若干。}
\da{答曰：得半徑一百二十步。}
\shu{術曰：以減從翻法開平方，置一於左上為上法，倍隅法為廉法，次商自之為隅法，併廉隅共為下法，與上法相乘除實。}
\shi{釋曰：此減從翻法開平方也。以減從翻法開平方，得圓徑。其法以初商一百為上法，倍隅法十二萬五千為廉法，得二十五萬。約次商得二十，以次商二十乘隅算得一萬，倍之得二萬五千為隅法。併廉法二十五萬、隅法二萬五千，共二十七萬五千為下法。以下法二十七萬五千與上法二十相乘，得五百五十萬，除實盡。故得半徑一百二十步，倍之得全徑二百四十步也。}
\end{problem}

\begin{problem}[\ 減從翻法開平方第二問]
\wen{二十與上次法相乘，得四千四百為益實。添入餘積，共五千六百為實。置一併廉法，從方共二百八十為下法，與上次法相乘除實盡。問半徑若干。}
\da{答曰：得半徑一百二十步。}
\shu{術曰：以減從翻法開平方，從方添入益實，併廉法為下法，與上法相乘除實。}
\shi{釋曰：此減從翻法開平方也。以減從翻法開平方，得圓徑。其術以二十與上次法二百二十相乘，得四千四百為益實。添入餘積一千二百，共五千六百為實。置一併廉法二百八十為下法，與上次法二十相乘，得五千六百，除實盡。故得半徑一百二十步，倍之得全徑二百四十步也。減從翻法之術，以益實添入積數，然後開方，此其所以為變術也。}
\end{problem}

\clearpage

%% ========================================================================
%% Chapter 5: Gou gu xiang cheng beizhi fa
%% ========================================================================
\subsection{勾股相乘倍之為實開平方}
\noindent
\textbf{總論：}勾股相乘倍之為實者，以勾股相乘之數倍之為實，開平方得弦，再以弦求容圓也。

此術與通勾股求容圓之術相類，而其步驟則異。通勾股求容圓，直以勾股相乘倍之為實，勾股和為法，除之得圓徑。此術則先開平方得弦，再以弦減勾股和得圓徑。二術雖異，其理則同，皆弦和較為容圓徑也。

勾股相乘倍之為實開平方之術，在測圓海鏡中為基本之法。學者先明此術，然後可學減從翻法、帶從開方等變術也。蓋變術生於正術，正術既明，變術可迎刃而解矣。

\bigskip

\begin{problem}[\ 勾股相乘倍之開平方第一問]
\wen{甲出南門東行一百二十步，乙出東門南行八十步，相望見之。問城徑。}
\da{答曰：城徑二百四十步。}
\shu{術曰：勾股相乘倍之為實，開平方得弦，併勾股為弦和，弦和較即容圓徑。}
\shi{釋曰：此勾股相乘倍之為實開平方也。以勾股相乘倍之為實，開平方得弦，再以弦求容圓。甲東行一百二十步為勾，乙南行八十步為股。以勾股相乘得九千六百，倍之得一萬九千二百為實。開平方得弦約一百三十八步半有奇。以勾股和二百減弦，餘約六十一步半有奇，為容圓徑。以此比例推之，得城全徑二百四十步也。}
\end{problem}

\begin{problem}[\ 勾股相乘倍之開平方第二問]
\wen{甲出北門東行一百五十步，乙出東門北行一百步，相望見之。問城徑。}
\da{答曰：城徑二百四十步。}
\shu{術曰：勾股相乘倍之為實，開平方得弦，併勾股為弦和，弦和較即容圓徑。}
\shi{釋曰：此勾股相乘倍之為實開平方也。以勾股相乘倍之為實，開平方得弦，再以弦求容圓。甲東行一百五十步為勾，乙北行一百步為股。以勾股相乘得一萬五千，倍之得三萬為實。開平方得弦約一百七十二步半有奇。以勾股和二百五十減弦，餘約七十七步半有奇，為容圓徑。以此比例推之，得城全徑二百四十步也。}
\end{problem}

\clearpage

%% ========================================================================
%% Chapter 6: Xian he jiao qiu yuanjing
%% ========================================================================
\subsection{弦和較求圓徑}
\noindent
\textbf{總論：}弦和較求圓徑者，以弦與勾股和之差，即容圓徑也。

弦和較者，何謂也？弦者，勾股弦也；和者，勾股和也；較者，差也。弦和較者，弦減勾股和之餘數也。以弦減勾股和，其餘即容圓徑，此測圓海鏡之要術也。

何以知弦和較即容圓徑也？蓋以勾股容圓之定理知之。勾股之內容圓，其直徑等於弦與勾股和之差，此定理也。以通勾股證之：通勾三百二十，通股六百，通弦六百八十。以通弦六百八十減通勾股和九百二十，餘二百四十，即圓城全徑。以此定理推之，無論何種勾股，其弦和較必為容圓徑，此千古不易之理也。

\bigskip

\begin{problem}[\ 弦和較求圓徑第一問]
\wen{甲出南門直行，乙出東門直行，相望見之。問以弦和較求圓徑。}
\da{答曰：城徑二百四十步。}
\shu{術曰：弦和較即容圓徑。弦者勾股弦也，和者勾股和也。弦減勾股和之較即圓徑。}
\shi{釋曰：此弦和較求圓徑也。弦和較者，通弦減通勾股和之較也，即容圓徑。其理至明，無待繁言。弦和較之為術，在測圓海鏡中最為簡要，以其一直可得，不須繁複計算也。然其理至深，非精通勾股者不能明其所以然也。}
\end{problem}

\begin{problem}[\ 弦和較求圓徑第二問]
\wen{甲出北門直行，乙出西門直行，相望見之。問以弦和較求圓徑。}
\da{答曰：城徑二百四十步。}
\shu{術曰：弦和較即容圓徑。弦者勾股弦也，和者勾股和也。弦減勾股和之較即圓徑。}
\shi{釋曰：此弦和較求圓徑也。以弦和較之數，即容圓徑。弦和較之術，為測圓之要術。以弦減勾股和，其餘即圓徑，此定理也。無論何種勾股，皆準此定理求之，無有不合者。}
\end{problem}

\clearpage

%% ========================================================================
%% Reference Table for Volume 2
%% ========================================================================
\subsection{卷二名數表}
\noindent
\textbf{卷二所用名數對照表：}

\begin{center}
\begin{longtable}{|l|l|l|}
\hline
\textbf{名稱} & \textbf{數值} & \textbf{說明} \\
\hline
\endfirsthead
\hline
\textbf{名稱} & \textbf{數值} & \textbf{說明} \\
\hline
\endhead
通股 & 六百步 & 通股之數 \\
\hline
邊股 & 三百步 & 邊股之數，通股之半 \\
\hline
底股 & 四百八十步 & 底股之數 \\
\hline
黃廣股 & 四百五十步 & 黃廣股之數 \\
\hline
黃長股 & 一百二十步 & 黃長股之數 \\
\hline
高股 & 二百步 & 高股之數 \\
\hline
平股 & 八十步 & 平股之數 \\
\hline
大差股 & 一百六十步 & 大差股之數 \\
\hline
小差股 & 六十步 & 小差股之數 \\
\hline
皇極股 & 二百四十步 & 皇極股之數 \\
\hline
太虛股 & 三十步 & 太虛股之數 \\
\hline
明股 & 八十步 & 明股之數 \\
\hline
重股 & 十六步 & 重股之數 \\
\hline
通弦 & 六百八十步 & 通弦之數 \\
\hline
邊弦 & 三百四十步 & 邊弦之數 \\
\hline
底弦 & 五百一十步 & 底弦之數 \\
\hline
黃廣弦 & 五百一十步 & 黃廣弦之數 \\
\hline
黃長弦 & 二百零八步 & 黃長弦之數 \\
\hline
\end{longtable}
\end{center}

\clearpage


%% ========================================================================
%% VOLUME 3 (卷三)
%% ========================================================================
\section{卷三：帶從開方及三乘方諸法}
\markboth{卷三}{卷三}

\noindent
\textbf{卷三總論：}卷三共分七章，專論\textbf{帶從開平方法}及\textbf{三乘方}之法，凡二十餘題。帶從開平方法者，以實較大於法，用帶從之法，開平方得圓城之徑也。又論減從翻法、益積、減積諸法，皆開方之變術也。

卷三之術，較卷二為更深。卷二論兩股兩弦及減從翻法開平方，術已較深。卷三則論帶從開平方、三乘方、益積減積、從方從廉隅法等，其理更深，其術更密。帶從開平方者，實內帶有從方之數，開平方時須帶從方一同開之，故名帶從開平方也。

三乘方者，三次方程之開方也。一乘方為平方，二乘方為立方，三乘方為立方以上之開方也。三乘方之術，以三次方程式求未知數，其理至精，其術至密。此天元術之精髓也，學者所當潛心玩索者也。

\bigskip

%% ========================================================================
%% Chapter 1: Dai cong kaipingfang fa
%% ========================================================================
\subsection{帶從開平方法}
\noindent
\textbf{總論：}帶從開平方法者，以帶從之實，用開平方之法，得圓城之徑也。所謂帶從者，實內帶有從方之數也。

帶從開平方之術，較正術為難。正術以實開平方，不帶從方。帶從開平方則實內帶有從方之數，開平方時須將從方一併開之。其法以實為主，從方為從，開平方時以從方輔助實數，漸次求得平方之根。

帶從開平方之步驟：一曰立實，二曰立從方，三曰初商，四曰置上法，五曰以次商乘從方為從法，六曰併實從共為總實，七曰以下法與上法相乘除總實。實盡則得半徑，不盡則繼續開之。此術較減從翻法更進一步，為天元術之深法也。

\bigskip

\begin{problem}[\ 帶從開平方法第一問]
\wen{乙出南門直行一百三十五步，甲出北門東行二百步見之。問以帶從開平方求城徑。}
\da{答曰：城徑二百四十步。}
\shu{術曰：為實以南行三百六十為從方，作帶從開平方法除之得全徑。}
\shi{釋曰：此帶從開平方法也。以從方之數，開平方得圓徑。其法以南行一百三十五步為實，以東行二百步與南行一百三十五步之和三百六十為從方。作帶從開平方法，以從方輔助實數，開平方得全徑二百四十步。帶從開平方之術，以從方為輔助，其理甚深，學者宜細心體會之。}
\end{problem}

\begin{problem}[\ 帶從開平方法第二問]
\wen{甲出北門東行一百二十步，乙出東門南行一百步，相望見之。問以帶從開平方求城徑。}
\da{答曰：城徑二百四十步。}
\shu{術曰：為實以兩行之和為從方，作帶從開平方法除之得全徑。}
\shi{釋曰：此帶從開平方法也。以從方之數，開平方得圓徑。甲東行一百二十步，乙南行一百步，兩行之和二百二十步為從方。以帶從開平方法開之，得全徑二百四十步。帶從開平方之術，雖較複雜，而其理則與正術相通。正術既明，帶從之術可迎刃而解也。}
\end{problem}

\begin{problem}[\ 帶從開平方法第三問]
\wen{甲出西門南行二百步，乙出南門西行一百五十步，相望見之。問以帶從開平方求城徑。}
\da{答曰：城徑二百四十步。}
\shu{術曰：為實以兩行之和為從方，作帶從開平方法除之得全徑。}
\shi{釋曰：此帶從開平方法也。以從方之數，開平方得圓徑。甲南行二百步，乙西行一百五十步，兩行之和三百五十步為從方。以帶從開平方法開之，得全徑二百四十步。帶從開平方之術，在測圓海鏡中為重要之法，以其能解較複雜之問題也。}
\end{problem}

\begin{problem}[\ 帶從開平方法第四問]
\wen{甲出南門直行二百五十步，乙出東門直行二百步，相望見之。問以帶從開平方求城徑。}
\da{答曰：城徑二百四十步。}
\shu{術曰：為實以兩行之和為從方，作帶從開平方法除之得全徑。}
\shi{釋曰：此帶從開平方法也。以從方之數，開平方得圓徑。甲南行二百五十步，乙東行二百步，兩行之和四百五十步為從方。以帶從開平方法開之，得全徑二百四十步。帶從開平方之術，以兩行之和為從方，其理在於以和輔實，以實求根，此開方之深術也。}
\end{problem}

\clearpage

%% ========================================================================
%% Chapter 2: Jian cong fanfa kaipingfang
%% ========================================================================
\subsection{減從翻法開平方}
\noindent
\textbf{總論：}減從翻法開平方者，以實內減去從方之數，翻轉開平方，得圓城之徑也。凡言減從翻法開平方者俱倣此。

減從翻法之術，在卷二中已略言之，此卷則更詳其法。減從翻法者，以實之數內減去從方之數，然後翻轉開平方。其名曰減從者，以從方減實也；名曰翻法者，以正術翻轉也。此術為開方之變術，較正術為難，學者須細心體會。

減從翻法之義：開平方之正術，以商除實，實不減商。減從翻法則以從方減實，實減從方而後開之。其所以須減者，以實之數過大，不減則不能開也。其所以須翻者，以減從之後，正術不可用，故須翻轉其法也。此變術之所以為變也。

\bigskip

\begin{problem}[\ 減從翻法開平方第一問]
\wen{十餘一千五百二十為負積。又以初商一百反減餘，從四十四餘五十六為負。從次商二十併負，從共七十六為下法，亦通後凡言減從翻法開平方者俱倣此。問半徑若干。}
\da{答曰：得半徑一百二十步。}
\shu{術曰：以減從翻法開平方曰：初商一百，置一於左上為上法。置一乘隅法，得四為隅法，作負隅開平方。}
\shi{釋曰：此減從翻法開平方也。以減從翻法開平方，得圓徑。其法以初商一百為上法，置一於左上。以隅法四乘之得四為隅法。以初商一百反減從方四十四，餘五十六為負從。以次商二十併負從五十六，共七十六為下法。以下法七十六與上法二十相乘，得一千五百二十，為負積。故得半徑一百二十步，倍之得全徑二百四十步也。}
\end{problem}

\begin{problem}[\ 減從翻法開平方第二問]
\wen{餘實五百五十萬，倍隅法得二十五萬為廉法。約次商得二十，置一於左次為上法，置一乘隅得二萬五千，併廉法共二十七萬五千為下法，與上法相乘除實。盡後如此類者倣此。問半徑若干。}
\da{答曰：得半徑一百二十步。}
\shu{術曰：以減從翻法開平方，置一於左上為上法，倍隅法為廉法，次商自之為隅法，併廉隅共為下法，與上法相乘除實。}
\shi{釋曰：此減從翻法開平方也。以減從翻法開平方，得圓徑。其法以餘實五百五十萬為實，倍隅法十二萬五千得二十五萬為廉法。約次商得二十，置一於左次為上法。以次商二十乘隅法得一萬，倍之得二萬五千為隅法。併廉法二十五萬、隅法二萬五千，共二十七萬五千為下法。以下法二十七萬五千與上法二十相乘，得五百五十萬，除實盡。故得半徑一百二十步，倍之得全徑二百四十步。}
\end{problem}

\begin{problem}[\ 減從翻法開平方第三問]
\wen{二十與上次法相乘，得四千四百為益實。添入餘積共五千六百為實。置一併廉法從方共二百八十為下法，與上次法相乘除實盡。問半徑若干。}
\da{答曰：得半徑一百二十步。}
\shu{術曰：以減從翻法開平方，從方添入益實，併廉法為下法，與上法相乘除實。}
\shi{釋曰：此減從翻法開平方也。以減從翻法開平方，得圓徑。其法以次商二十與上次法二百二十相乘，得四千四百為益實。添入餘積一千二百，共五千六百為實。置一併廉法從方共二百八十為下法。以下法二百八十與上次法二十相乘，得五千六百，除實盡。故得半徑一百二十步，倍之得全徑二百四十步。減從翻法之術，以益實添入積數，然後開方，此其所以為變術也。}
\end{problem}

\clearpage

%% ========================================================================
%% Chapter 3: San cheng fang fa
%% ========================================================================
\subsection{三乘方法}
\noindent
\textbf{總論：}三乘方法者，以三次方程之法，求圓城之徑也。所謂三乘者，立方以上之開方也。

三乘方者，何謂也？一乘方曰平方，二乘方曰立方，三乘方曰三乘。三乘方之術，以三次方程式求未知數。其方程式之形為：實 = 從方 × 未知數 + 廉法 × 未知數² + 隅法 × 未知數³。以開方之法，漸次求得未知數之值，即圓城之半徑也。

三乘方之術，在測圓海鏡中為最深之法。其術須用從方、廉法、隅法三者相配，以開立方以上之方。從方者，一次項之係數也；廉法者，二次項之係數也；隅法者，三次項之係數也。三者相配，以開三乘方，其理至深，其術至密，非精通天元術者不能明也。

\bigskip

\begin{problem}[\ 三乘方法第一問]
\wen{置一自之以乘，從二廉，置一自乘再乘為隅法。併從方廉隅共七千九十八萬一千三百四十為下法，與上法相乘除實。問半徑若干。}
\da{答曰：得半徑一百二十步。}
\shu{術曰：以三乘方法開之，置一自之為上法，以次商乘廉法為從方，再以次商自之為隅法，併從方廉隅為下法，與上法相乘除實。}
\shi{釋曰：此三乘方法也。以三乘方法開之，得圓徑。其法置一自之為上法，以從二廉為從方，置一自乘再乘為隅法。併從方、廉法、隅法共七千九十八萬一千三百四十為下法。與上法相乘除實，得半徑一百二十步，倍之得全徑二百四十步。三乘方之術，以從方、廉法、隅法三者相配，其理至深，此天元術之精髓也。}
\end{problem}

\begin{problem}[\ 三乘方法第二問]
\wen{置一自之以乘從二廉，置一自乘再乘為隅法。併從方廉隅共七千九十八萬一千三百四十為下法，與上法相乘除實盡。問半徑若干。}
\da{答曰：得半徑一百二十步。}
\shu{術曰：以三乘方法開之，置一自之為上法，以次商乘廉法為從方，再以次商自之為隅法，併從方廉隅為下法，與上法相乘除實。}
\shi{釋曰：此三乘方法也。以三乘方法開之，得圓徑。其術與上問同，特其實數異耳。三乘方之術，以三次方程求未知數，其理至精。置一自之為上法，次商乘廉法為從方，次商自之為隅法，三者相配，以開三乘方。學者能明此術，則天元之學思過半矣。}
\end{problem}

\clearpage

%% ========================================================================
%% Chapter 4: Yi ji jian ji fa
%% ========================================================================
\subsection{益積減積法}
\noindent
\textbf{總論：}益積減積法者，以益積或減積之法，開方得圓城之徑也。益積者添入積數，減積者減去積數。

益積減積之術，為開方之變術。正術以實開方，不增不減。益積減積則或添入積數，或減去積數，然後開方。其名益積者，以積數益之使大也；名減積者，以積數減之使小也。或益或減，視實之大小而定，實小則益之，實大則減之。

益積減積之義：開平方之時，有時實之數過小，不足以開平方，則須益積以助之。有時實之數過大，開平方之後仍有餘積，則須減積以平之。或益或減，皆所以使實之數適合於開平方也。此變術之妙用也。

\bigskip

\begin{problem}[\ 益積減積法第一問]
\wen{二十與上次法相乘，得四千四百為益實。添入餘積，共五千六百為實。置一併廉法，從方共二百八十為下法，與上次法相乘除實盡。問半徑若干。}
\da{答曰：得半徑一百二十步。}
\shu{術曰：以益積法添入益實，併廉法為下法，與上法相乘除實。}
\shi{釋曰：此益積法也。以益積之數，開平方得圓徑。其法以次商二十與上次法二百二十相乘，得四千四百為益實。添入餘積一千二百，共五千六百為實。置一併廉法從方共二百八十為下法。以下法二百八十與上次法二十相乘，得五千六百，除實盡。故得半徑一百二十步，倍之得全徑二百四十步。益積之法，以益實添入積數，使實之數足供開方，此變術之妙也。}
\end{problem}

\begin{problem}[\ 益積減積法第二問]
\wen{餘實五百五十萬，倍隅法得二十五萬為廉法。約次商得二十，置一於左次為上法，置一乘隅得二萬五千，併廉法共二十七萬五千為下法，與上法相乘除實。問半徑若干。}
\da{答曰：得半徑一百二十步。}
\shu{術曰：以減積法減去積數，併廉法為下法，與上法相乘除實。}
\shi{釋曰：此減積法也。以減積之數，開平方得圓徑。其法以餘實五百五十萬為實，倍隅法得二十五萬為廉法。約次商得二十，置一於左次為上法。以次商二十乘隅法得一萬，倍之得二萬五千為隅法。併廉法二十五萬、隅法二萬五千，共二十七萬五千為下法。以下法二十七萬五千與上法二十相乘，得五百五十萬，除實盡。故得半徑一百二十步，倍之得全徑二百四十步。減積之法，以減積使實之數適合於開方，此亦變術之妙也。}
\end{problem}

\clearpage

%% ========================================================================
%% Chapter 5: Cong fang cong lian yu fa
%% ========================================================================
\subsection{從方從廉隅法}
\noindent
\textbf{總論：}從方從廉隅法者，以從方、從廉、隅法三者之關係，開方得圓城之徑也。

從方從廉隅之術，為三乘方之基礎。從方者，一次項也；從廉者，二次項也；隅法者，三次項也。三者相配，以成三次方程式，然後開三乘方，求得未知數。此術在測圓海鏡中為最深之法，學者須潛心玩索，方能明其奧妙。

從方從廉隅之義：開平方之時，有從方無從廉無隅法。開立方之時，有從方有從廉無隅法。開三乘方之時，有從方有從廉有隅法。三者漸次增加，其術漸次複雜。學者須先明開平方，次明開立方，然後可學開三乘方也。此學術之次序，不可躐等也。

\bigskip

\begin{problem}[\ 從方從廉隅法第一問]
\wen{置一自之以乘從二廉，置一自乘再乘為隅法。併從方廉隅共七千九十八萬一千三百四十為下法，與上法相乘除實盡。問半徑若干。}
\da{答曰：得半徑一百二十步。}
\shu{術曰：以從方從廉隅法開之，置一自之為上法，以次商乘廉法為從方，再以次商自之為隅法，併從方廉隅為下法，與上法相乘除實。}
\shi{釋曰：此從方從廉隅法也。以從方從廉隅法開之，得圓徑。其法置一自之為上法，以從二廉為從方，置一自乘再乘為隅法。併從方、廉法、隅法共七千九十八萬一千三百四十為下法。與上法相乘除實，得半徑一百二十步，倍之得全徑二百四十步。從方從廉隅之術，三者相配，其理至深，此天元術之最高深處也。}
\end{problem}

\begin{problem}[\ 從方從廉隅法第二問]
\wen{置一自之為上法，置一乘隅得二萬五千，併廉法共二十七萬五千為下法，與上法相乘除實。問半徑若干。}
\da{答曰：得半徑一百二十步。}
\shu{術曰：以從方從廉隅法開之，置一自之為上法，以次商乘廉法為從方，再以次商自之為隅法，併從方廉隅為下法，與上法相乘除實。}
\shi{釋曰：此從方從廉隅法也。以從方從廉隅法開之，得圓徑。其術與上問相類，特其數異耳。從方從廉隅之術，雖為最深，而其理則與開平方相通。開平方既明，開三乘方可迎刃而解，特其步驟較繁耳。學者須耐心體會，不可畏難而退也。}
\end{problem}

\clearpage

%% ========================================================================
%% Chapter 6: Bu ji di gou bian gu tiao
%% ========================================================================
\subsection{不及底勾邊股條}
\noindent
\textbf{總論：}不及底勾邊股條者，以底勾、邊股不及之數，用測望之法，求圓城之徑也。

不及者，何謂也？不及者，不到也，未達也。底勾邊股之不及者，以底勾之數不及邊股之數，或邊股之數不及底勾之數，兩數相差，用測望之法，求得圓城之徑。此術較為複雜，須用測望之法，以兩數之差推算圓徑。

底勾邊股條之術，在測圓海鏡中為較深之法。其術以底勾之數與邊股之數相比較，以其不及之數，用測望之法推算圓徑。測望之法，以目測之距離推算實際之距離，此古代測量之要術也。學者須先明測望之理，然後可學此術也。

\bigskip

\begin{problem}[\ 不及底勾邊股條第一問]
\wen{出西門南行二百二十五步，有塔出北門東行六十四步望塔正居城之半。問城徑。}
\da{答曰：城徑二百四十步。}
\shu{術曰：此以不及底勾與不及邊股測望也。南行二百二十五步與高股同，即半徑為勾之股。東行六十四步與平勾同，即半徑為股之勾也。當以平勾高股立法。}
\shi{釋曰：此以不及底勾邊股條求城徑也。以底勾邊股之不及數，測望得圓徑。出西門南行二百二十五步為高股，出北門東行六十四步為平勾。高股者，半徑為勾之股也；平勾者，半徑為股之勾也。以平勾高股測望，得圓徑二百四十步。此測望之法，以目測推算實距，其理至精。}
\end{problem}

\begin{problem}[\ 不及底勾邊股條第二問]
\wen{乙從城外西南坤隅南行三百六十步，甲出北門東行二百步見之。問城徑。}
\da{答曰：城徑二百四十步。}
\shu{術曰：二行相乘即半徑冪。}
\shi{釋曰：此以底勾大差股立法測望也。乙從坤隅南行大差股也，甲東行底勾也。底勾為城北半勾，大差股為城西南虛股。以二行相乘即半徑冪，開平方得半徑一百二十步，倍之得全徑二百四十步。此術以二行相乘得半徑冪，開平方得半徑，其理至妙。}
\end{problem}

\begin{problem}[\ 不及底勾邊股條第三問]
\wen{甲出北門東行底勾也。底勾為城北半勾，明股為城南餘股。問城徑。}
\da{答曰：城徑二百四十步。}
\shu{術曰：以底勾明股之數，約之得圓徑。}
\shi{釋曰：此底勾明股立法也。以底勾明股之數，測望得圓徑。底勾為城北半勾，明股為城南餘股。以底勾明股之數，約之得圓徑二百四十步。此不及底勾邊股條之術，以兩數之差推算圓徑，其理至深。}
\end{problem}

\clearpage

%% ========================================================================
%% Chapter 7: Shesheng buta chengfa
%% ========================================================================
\subsection{測望不迭乘法}
\noindent
\textbf{總論：}測望不迭乘法者，以測望之法，不迭次相乘，求圓城之徑也。

不迭乘者，何謂也？不迭乘者，不連續相乘也。測望之時，有須連續相乘者，有不必連續相乘者。此術以不迭乘法，分別從方從廉，開平方得圓徑。其術較簡，而其理則與迭乘法無異。

測望不迭乘法之義：測望之法，以目測推算實際距離。有時須連續相乘，有時不須連續相乘。不迭乘法者，以一次乘法得之，不須反覆迭乘。其術較簡便，而其理則不減。學者能明此術，則測望之學思過半矣。

\bigskip

\begin{problem}[\ 測望不迭乘法第一問]
\wen{此法分別從方從廉明白，故重錄附之。問以測望不迭乘法求半徑若干。}
\da{答曰：得半徑一百二十步。}
\shu{術曰：以測望之法，分別從方從廉，不迭次相乘，開平方得圓徑。}
\shi{釋曰：此測望不迭乘法也。以測望之法，不迭次相乘，開平方得圓徑。其法分別從方從廉，不以迭乘，而以一次乘法得之。其術較簡，而其理則與迭乘法無異。此術之所以為善也，以其簡便易行，而不失其精確也。}
\end{problem}

\begin{problem}[\ 測望不迭乘法第二問]
\wen{甲出南門東行一百步，乙出東門南行六十步，相望見之。問以測望不迭乘法求城徑。}
\da{答曰：城徑二百四十步。}
\shu{術曰：以測望之法，分別從方從廉，不迭次相乘，開平方得圓徑。}
\shi{釋曰：此測望不迭乘法也。以測望之法，不迭次相乘，開平方得圓徑。甲東行一百步為勾，乙南行六十步為股。以測望之法，分別從方從廉，不迭次相乘，開平方得圓徑二百四十步。此術簡便易行，為測望之法中較善者也。}
\end{problem}

\clearpage

%% ========================================================================
%% Reference Table for Volume 3
%% ========================================================================
\subsection{卷三名數表}
\noindent
\textbf{卷三所用名數對照表：}

\begin{center}
\begin{longtable}{|l|l|l|}
\hline
\textbf{名稱} & \textbf{數值} & \textbf{說明} \\
\hline
\endfirsthead
\hline
\textbf{名稱} & \textbf{數值} & \textbf{說明} \\
\hline
\endhead
上法 & 初商自之 & 開方之上位法 \\
\hline
從方 & 次商乘廉法 & 從方之數 \\
\hline
廉法 & 倍隅法 & 廉法之數 \\
\hline
隅法 & 次商自之 & 隅法之數 \\
\hline
下法 & 從方廉隅共 & 下法之數 \\
\hline
實 & 積數 & 實之數 \\
\hline
益實 & 添入之數 & 益實之數 \\
\hline
減實 & 減去之數 & 減實之數 \\
\hline
負積 & 負數之積 & 負積之數 \\
\hline
帶從 & 從方帶入 & 帶從之數 \\
\hline
平方 & 二次方 & 開平方 \\
\hline
立方 & 三次方 & 開立方 \\
\hline
三乘方 & 四次方 & 開三乘方 \\
\hline
初商 & 第一次商數 & 開方之初步 \\
\hline
次商 & 第二次商數 & 開方之次步 \\
\hline
半徑 & 一百二十步 & 圓城之半徑 \\
\hline
全徑 & 二百四十步 & 圓城之全徑 \\
\hline
\end{longtable}
\end{center}

\clearpage

%% ========================================================================
%% END MATTER
%% ========================================================================
\section*{跋}
\addcontentsline{toc}{section}{跋（Postscript）}
\markboth{跋}{跋}

測圓海鏡分類釋術十卷，元李冶撰，明顧應祥分類釋術。其書以勾股測圓之術，分類編次，以便學者。卷一至卷三論通勾股、邊勾股、底股弦、皇極勾股、太陰勾股、明勾股、重勾股諸法求容圓，及兩股求容圓、兩弦求容圓、帶從開平方、減從翻法開平方、三乘方法、益積減積法、從方從廉隅法諸術。

其術之大要，以勾股相乘倍之為實，勾股和為法，除之得圓徑。又以弦和較即容圓徑。其開方之法，有帶從、減從翻法、益積、減積諸術，皆天元術之支流也。

顧應祥序謂：李冶原書每條下細草，俱徑立天元一，反覆合之，而無下手之術。使後學之士茫然無門路可入。故去其細草，立一算術，又以其所立通勾邊股之屬各以類分之。此誠後學之津梁也。

夫測圓之術，源遠流長。自《周髀算經》載勾股之術，歷代相傳，代有發明。漢劉徽注《九章》，創割圓之術。魏趙爽釋《周髀》，明弦圖之義。宋秦九韶著《數書九章》，立正負開方之術。元李冶著《測圓海鏡》，以天元術測圓，集大成焉。明顧應祥分類釋術，使奧義顯明，功不可沒。

測圓海鏡之學，以勾股為體，以容圓為用。其勾股之類，有通勾股、邊勾股、底勾股、黃廣勾股、黃長勾股、高勾股、平勾股、大差勾股、小差勾股、皇極勾股、太虛勾股、明勾股、重勾股等十餘種。各以其所在之位而得名，其理則一以貫之，無非勾股而已。

其容圓之術，以弦和較為圓徑。弦和較者，弦與勾股和之差也。無論何種勾股，其弦和較必為容圓徑，此千古不易之理也。以通勾股證之：通勾三百二十，通股六百，通弦六百八十。以通弦六百八十減通勾股和九百二十，餘二百四十，即圓城全徑。以此推之，百無一失。

顧應祥之分類釋術，以術為綱，以題為目，使學者由術以求題，不由題以求術。其分類之法，以通勾股求容圓為一類，邊勾股求容圓為一類，底股弦求通圓徑為一類，皇極勾股求容圓為一類，太陰勾股求容圓為一類，明勾股求容圓為一類，重勾股求容圓為一類。每類之前，各立總論，詳述其術之理。每題之下，分問、答、術、釋四段，使學者一目了然。此誠善於教學者也。

\bigskip
\begin{flushright}
測圓海鏡分類釋術卷三終
\end{flushright}

\clearpage

%% ========================================================================
%% COLOPHON
%% ========================================================================
\thispagestyle{empty}
\begin{center}
\vspace*{3cm}
{\Large 測圓海鏡分類釋術}

\vspace{1cm}
{\large 卷一至卷三終}

\vspace{2cm}
\shortline

\vspace{1cm}
{\normalsize
\begin{tabular}{rl}
總校官候補中書 & 臣吳紹潔 \\
校對官中官正 & 臣郭長發 \\
謄錄監生 & 臣丁湘錦 \\
\end{tabular}
}

\vspace{2cm}
{\small\textcolor{notecolor}{%
此為欽定四庫全書薈要本\\
子部天文算法類\\
清乾隆年間鈔本
}}

\end{center}

\clearpage

%% ========================================================================
%% MODERN EDITORIAL NOTE
%% ========================================================================
\section*{Modern Editorial Note}
\addcontentsline{toc}{section}{Modern Editorial Note}
\markboth{Modern Editorial Note}{Modern Editorial Note}

This edition presents the first three volumes（卷一至卷三）of the \textit{Ceyuan Haijing Fenlei Shishu}（測圓海鏡分類釋術）, a classified rearrangement of Li Ye's famous \textit{Ceyuan Haijing}（Sea Mirror of Circle Measurements）, prepared by Gu Yingxiang（顧應祥）during the Ming dynasty.

\subsection*{Historical Context}
The original \textit{Ceyuan Haijing} was composed by Li Ye（李冶, 1192--1279）in 1248 during the Yuan dynasty. It is one of the masterpieces of medieval Chinese algebra, presenting 170 problems concerning the computation of the diameter of an inscribed circle in a right triangle and its numerous variations. Li Ye's original work employed the \textit{tian yuan shu}（天元術, ``celestial element method''）--- an advanced algebraic technique using rod numerals to set up and solve polynomial equations.

\subsection*{Gu Yingxiang's Rearrangement}
Gu Yingxiang（顧應祥, 1483--1565）, a Ming dynasty mathematician and Minister of Justice, found Li Ye's original presentation difficult for students to follow because the \textit{tian yuan shu} method had been largely forgotten. He therefore:
\begin{enumerate}
\item Removed the detailed ``fine grass''（細草）working using \textit{tian yuan shu}
\item Replaced it with explicit step-by-step arithmetic procedures（算術）
\item Reorganized the 170 problems by mathematical technique rather than geometric configuration
\item Added explanatory notes（釋術）explaining why each method works
\end{enumerate}

\subsection*{Structure of Volumes 1--3}
\begin{itemize}
\item \textbf{Volume 1} focuses on ``containing-circle''（容圓）methods using various types of right triangles: \textit{tong}（通）, \textit{bian}（邊）, \textit{di}（底）, \textit{huangji}（皇極）, \textit{taiyin}（太陰）, \textit{ming}（明）, and \textit{chong}（重）gougu configurations.
\item \textbf{Volume 2} covers methods using two legs（兩股）, two hypotenuses（兩弦）, difference of legs, and the ``subtractive-associate inverted method''（減從翻法）for root extraction.
\item \textbf{Volume 3} presents advanced root-extraction techniques: the ``accompanying-associate''（帶從）method, three-fold multiplication（三乘方）, additive and subtractive accumulation（益積減積）, and methods involving the edge-leg and base-leg relationships.
\end{itemize}

\subsection*{Mathematical Significance}
The problems all revolve around a single geometric configuration: a circular city（圓城）with given measurements from various observation points, from which one must determine the city's diameter. The invariant answer throughout is \textbf{240 bu}（步）, demonstrating the algebraic relationships rather than varying geometric scenarios.

The methods exemplify the sophisticated Chinese tradition of root extraction（開方術）and polynomial algebra that predated European developments by centuries. Gu Yingxiang's classification system reveals the underlying structural unity of Li Ye's 170 problems, organizing them by their mathematical technique rather than their surface geometric appearance.

\subsection*{Technical Note on Typesetting}
This edition was typeset using XeLaTeX with Noto Sans CJK SC font support. The original Siku Quanshu Huiyao manuscript employs traditional Chinese characters; this edition uses simplified Chinese characters（簡體字）for broader accessibility while preserving all technical terminology in their standard mathematical forms.

\vspace{1cm}
\begin{flushright}
\textit{Typeset with XeLaTeX using Noto Sans CJK SC}\\
\textit{Modern edition prepared from the Siku Quanshu Huiyao manuscript}
\end{flushright}

\end{document}
